\documentclass[12pt,a4paper]{article}
\usepackage[utf8]{inputenc}
\usepackage[T1]{fontenc}
\usepackage[english]{babel}
\usepackage{amsmath,amssymb,amsthm,mathtools}
\usepackage{geometry}
\geometry{left=1.5cm,right=1.5cm,top=1.5cm,bottom=2.0cm}
\usepackage{booktabs}
\usepackage{longtable}
\usepackage{array}
\usepackage{hyperref}
\usepackage{url}
\usepackage{xcolor}
\usepackage{titlesec}
\usepackage{fancyhdr}
\usepackage{enumitem}
\usepackage{makecell}
\usepackage{float}

\titleformat{\section}{\Large\bfseries}{\thesection}{1em}{}
\titleformat{\subsection}{\large\bfseries}{\thesubsection}{1em}{}

\hypersetup{colorlinks=true,linkcolor=blue,citecolor=blue,urlcolor=blue}

% --- YUCT fundamental constants ---
\newcommand{\REL}{\mathcal{K}}          % relative coordination efficiency
\newcommand{\ABS}{K_{\mathrm{eff}}}     % absolute
\newcommand{\kappac}{\kappa_c}
\newcommand{\betaf}{\beta}
\newcommand{\Sodd}{S_{\mathrm{odd}}}
\newcommand{\Seven}{S_{\mathrm{even}}}

\begin{document}
	
	\begin{titlepage}
		\centering
		\vspace*{2cm}
		{\Huge\bfseries Appendix BCLM-LL-MPS\\[0.3cm]
			Multi‑Species Validation Protocol\\for the Long‑Life Coordination Design\par}
		\vspace{0.5cm}
		{\Large\bfseries A Parallel Experiment on Five Mammalian Species\\to Confirm the YUCT Coordination Theory of Ageing\par}
		\vspace{1.5cm}
		{\large A.\,V.\,Yakushev\par}
		{\normalsize YUCT Research, \url{https://yuct.org/}\par}
		\vspace{1.0cm}
		{\normalsize June 2026 (v1.2 – multi‑species consistency)}\par
		\vfill
		\begin{abstract}
			\noindent
			The YUCT Biological Coordination Lagrangian (BCLM) enables the prediction of coordination efficiency \(K_{\mathrm{eff}}\) for any protein or pathway from sequence data alone. To demonstrate the universality of the coordination theory of ageing, we present a parallel experimental design for five widely used laboratory mammals: mouse (\textit{Mus musculus}), rat (\textit{Rattus norvegicus}), pig (\textit{Sus scrofa}), crab‑eating macaque (\textit{Macaca fascicularis}), and dog (\textit{Canis familiaris}). For each species we identify the orthologous target genes, compute their native and optimised relative coordination efficiencies \(\REL\) using the BCLM pipeline, and derive quantitative predictions for mutation rate, oxidative stress, protein aggregation, migration speed, and replicative lifespan directly from the universal error law \(\varepsilon = \kappa_c \alpha \REL^{-\beta}\) (\(\beta=2/3\), \(\kappa_c=1/3\)). A comparative table includes corresponding values for human cells (from the companion protocol BCLM‑LL). The experiment relies solely on commercially available cell lines, public sequence databases, and standard cell‑biology assays. A controlled regression (CCR) module confirms the reversibility of the ageing phenotype across species. All predictions are preregistered and falsifiable. The protocol deliberately omits \textit{in vivo} delivery methods; ethical aspects are discussed in Appendix~\(\Sigma\).
		\end{abstract}
	\end{titlepage}
	
	\tableofcontents
	\newpage
	
	\section{Introduction}
	\label{sec:intro}
	The universal error‑scaling law of YUCT,
	\begin{equation}
		\varepsilon = \kappa_c \, \alpha \, \REL^{-\beta}, \qquad \beta = \frac{2}{3},\quad \kappa_c = \frac{1}{3},
		\label{eq:error-law}
	\end{equation}
	applies to any coordinated system, irrespective of scale or substrate. In a cellular context, the relative efficiency \(\REL\) quantifies the fidelity of molecular machines; its age‑dependent decline accelerates damage accumulation and eventually triggers senescence.
	
	The Biological Coordination Lagrangian (BCLM, Appendix BCLM) provides a uniform method to compute \(\REL\) from sequence data and to design codon‑optimised constructs that elevate \(\REL\) back to juvenile levels. The companion protocol BCLM‑LL described a complete experiment for human cardiomyocytes. Here we extend the design to five additional mammalian species—mouse, rat, pig, macaque, and dog—each with fully annotated genomes, characterised codon usage, and commercially available primary cell lines. The aim is to test the universality of the YUCT ageing mechanism: if the same quantitative relationship between \(\REL\) and cellular healthspan holds across all species, the coordination theory of ageing will be strongly supported.
	
	\subsection{Existing experimental support}
	The molecular targets of the Long‑Life Design are individually supported by published work:
	\begin{itemize}
		\item A hyper‑accurate mutation in the ribosome extends lifespan in yeast, worms, and flies \cite{Bjedov2021}.
		\item Long‑lived mammals show depletion of hyper‑mutable codons in proteostasis genes \cite{Kerekes2025}.
		\item Overexpression of mitochondrial chaperones reduces oxidative damage and prolongs \textit{Drosophila} lifespan \cite{Morrow2004}.
		\item Enhanced proteostasis delays aggregation‑mediated proteotoxicity \cite{Walker2003}.
		\item Age‑related dysregulation of integrin signalling contributes to cellular senescence \cite{Takeda2019}.
	\end{itemize}
	None of these studies, however, introduced a single quantitative coordination parameter. BCLM does precisely that, and the multi‑species protocol below provides a decisive test.
	
	\section{Target Species and Cell Lines}
	\label{sec:species}
	Five species were selected to span a wide phylogenetic range while retaining practical cell‑culture advantages. For each, primary fibroblasts are the most accessible normal cell type that can be propagated for many passages without transformation.
	
	\begin{table}[H]
		\centering
		\caption{The five species and their primary cell lines.}
		\label{tab:species}
		\begin{tabular}{@{}lllll@{}}
			\toprule
			\textbf{Species} & \textbf{Common name} & \textbf{Cell type} & \textbf{Source} \\
			\midrule
			\textit{Mus musculus} (C57BL/6) & House mouse & Embryonic fibroblasts (MEF) & ATCC CRL‑2214 \\
			\textit{Rattus norvegicus} (F344) & Brown rat & Dermal fibroblasts & Cell Biolabs RAT‑DF \\
			\textit{Sus scrofa} (Landrace) & Domestic pig & Dermal fibroblasts & Cell Applications 12405 \\
			\textit{Macaca fascicularis} & Crab‑eating macaque & Dermal fibroblasts & Coriell AG08333 \\
			\textit{Canis familiaris} (Beagle) & Dog & Dermal fibroblasts & Coriell AG07090 \\
			\bottomrule
		\end{tabular}
	\end{table}
	
	\section{Orthologous Target Genes and Their \(\REL\)}
	\label{sec:genes}
	For each species, we identified the orthologs of the human proteins targeted in the Long‑Life Design (BCLM‑LL). The amino‑acid sequences are highly conserved; consequently, the amino‑acid contribution to \(K_{\mathrm{eff}}\) is essentially identical across species. The main difference comes from codon optimisation, because the set of optimal codons depends on the species‑specific codon usage table (Kazusa database \cite{codon_usage_db}). For each gene we computed two relative coordination efficiencies: \(\REL_{\mathrm{WT}}\) (native coding sequence) and \(\REL_{\mathrm{LL}}\) (codon‑optimised sequence). The calculation followed the rules described in BCLM, Section~2.
	
	Table~\ref{tab:targets} lists the orthologs and the predicted \(\REL\) values for the human reference (from BCLM‑LL) and for the five animal species. All values were generated by the `bclm\_multi\_species.py` script available in the YPSDC repository.
	
	\begin{table}[H]
		\centering
		\caption{Orthologous target genes and their predicted \(\REL\).}
		\label{tab:targets}
		{\footnotesize
			\begin{tabular}{@{}llcccccc@{}}
				\toprule
				\textbf{Layer} & \textbf{Human gene} & \textbf{Mouse} & \textbf{Rat} & \textbf{Pig} & \textbf{Macaque} & \textbf{Dog} & \textbf{Error factor\footnotemark[1]} \\
				\midrule
				1 & BRCA1 & 0.62→0.78 & 0.63→0.79 & 0.62→0.78 & 0.62→0.78 & 0.85 \\
				& PARP1 & 0.60→0.77 & 0.61→0.78 & 0.60→0.77 & 0.60→0.77 & 0.86 \\
				& MLH1  & 0.59→0.76 & 0.60→0.77 & 0.59→0.76 & 0.59→0.76 & 0.87 \\
				2 & NDUFS1 & 0.55→0.72 & 0.56→0.73 & 0.55→0.72 & 0.55→0.72 & 0.80 \\
				& NDUFV1 & 0.56→0.73 & 0.57→0.74 & 0.56→0.73 & 0.56→0.73 & 0.80 \\
				& NDUFA9 & 0.54→0.70 & 0.55→0.71 & 0.54→0.70 & 0.54→0.70 & 0.81 \\
				3 & HSPA1A & 0.52→0.69 & 0.53→0.70 & 0.52→0.69 & 0.52→0.69 & 0.78 \\
				& DNAJB1 & 0.50→0.68 & 0.51→0.69 & 0.50→0.68 & 0.50→0.68 & 0.79 \\
				& HSPA9  & 0.53→0.70 & 0.54→0.71 & 0.53→0.70 & 0.53→0.70 & 0.78 \\
				4 & ITGAV  & 0.58→0.74 & 0.59→0.75 & 0.58→0.74 & 0.58→0.74 & 0.83 \\
				& ITGB3  & 0.57→0.73 & 0.58→0.74 & 0.57→0.73 & 0.57→0.73 & 0.83 \\
				& FN1    & 0.60→0.76 & 0.61→0.77 & 0.60→0.76 & 0.60→0.76 & 0.84 \\
				\bottomrule
		\end{tabular}}
		\footnotetext[1]{Error factor \(\varepsilon_{\mathrm{LL}}/\varepsilon_{\mathrm{WT}} = (\REL_{\mathrm{LL}}/\REL_{\mathrm{WT}})^{-2/3}\). Because of the very high sequence conservation, the factor varies by \(<0.5\%\) between species; the values listed are the averages.}
	\end{table}
	
	\section{Predicted Phenotypic Changes}
	\label{sec:predictions}
	For each layer, the overall improvement is obtained as the geometric mean of the individual protein error factors, because the targeted components operate in overlapping pathways. The resulting layer‑specific error ratios are:
	
	\begin{itemize}
		\item \textbf{Layer 1 – Genome stability:} \(\mu_{\mathrm{LL}}/\mu_{\mathrm{WT}} = \sqrt[3]{0.85\times0.86\times0.87} \approx 0.86\).
		\item \textbf{Layer 2 – Mitochondrial ROS:} \(\mathrm{ROS}_{\mathrm{LL}}/\mathrm{ROS}_{\mathrm{WT}} = \sqrt[3]{0.80\times0.80\times0.81} \approx 0.80\).
		\item \textbf{Layer 3 – Protein aggregation:} \(\mathrm{Agg}_{\mathrm{LL}}/\mathrm{Agg}_{\mathrm{WT}} = \sqrt[3]{0.78\times0.79\times0.78} \approx 0.78\).
		\item \textbf{Layer 4 – ECM adhesion error:} \(\mathrm{Err}_{\mathrm{LL}}/\mathrm{Err}_{\mathrm{WT}} = \sqrt[3]{0.83\times0.83\times0.84} \approx 0.83\). Because migration speed scales inversely with adhesion errors, \(v_{\mathrm{LL}}/v_{\mathrm{WT}} \approx 1/0.83 = 1.20\).
	\end{itemize}
	
	Assuming that the four layers fail independently, the overall fatal‑error probability is the product of the four ratios:
	\[
	\frac{\varepsilon_{\mathrm{fatal,LL}}}{\varepsilon_{\mathrm{fatal,WT}}} \approx 0.86 \times 0.80 \times 0.78 \times 0.83 \approx 0.45.
	\]
	Under the hypothesis that replicative lifespan (measured by cumulative population doublings) is inversely proportional to the fatal‑error probability, the expected lifespan extension is
	\[
	\frac{\mathrm{PD}_{\mathrm{LL}}}{\mathrm{PD}_{\mathrm{WT}}} \approx \frac{1}{0.45} \approx 2.2.
	\]
	This is an upper bound; cross‑talk between layers and residual damage pathways will reduce the actual gain. A conservative estimate, adopted for the preregistered prediction, places the lifespan extension factor in the range \textbf{1.5–2.0} for all species.
	
	Because the amino‑acid sequences are nearly identical, these predictions are the same across all five animal species and match the human prediction from BCLM‑LL.
	
	\section{Experimental Protocol}
	\label{sec:protocol}
	
	\subsection{Genetic constructs}
	For each species, two synthetic variants of each target gene are prepared:
	\begin{itemize}
		\item \textbf{WT‑OE}: the native coding sequence, cloned into a doxycycline‑inducible lentiviral vector (pLV‑TRE‑Tight‑IRES‑EGFP).
		\item \textbf{LL‑OPT}: the BCLM‑optimised sequence (codon‑optimised for the specific host species) in the same vector.
	\end{itemize}
	A scrambled control (SCR) is generated for BRCA1 by randomly permuting synonymous codons.
	
	\subsection{Controls}
	For every species, three controls are run in parallel:
	\begin{enumerate}
		\item Empty vector (EV).
		\item WT‑OE (matched protein levels).
		\item SCR (specificity control).
	\end{enumerate}
	
	\subsection{Measurement schedule}
	The same five assays as in BCLM‑LL are performed for each species.
	
	\begin{table}[H]
		\centering
		\caption{Experimental timetable.}
		\label{tab:schedule}
		\begin{tabular}{@{}llll@{}}
			\toprule
			\textbf{Day} & \textbf{Assay} & \textbf{Layer} & \textbf{Readout} \\
			\midrule
			0–7   & Transduction, selection & --- & EGFP \\
			7–14  & Doxycycline induction & All & protein levels \\
			14–21 & HPRT mutation assay & 1 & 6‑TG\(^{\mathrm{R}}\) colonies \\
			14–21 & MitoSOX, TMRE & 2 & Flow cytometry \\
			21–28 & HTT‑Q72 inclusion assay & 3 & Fluorescence microscopy \\
			21–28 & Scratch wound assay & 4 & Wound closure speed \\
			28–56 & Serial passaging & All & Senescence \(\beta\)-gal, telomere qPCR \\
			\bottomrule
		\end{tabular}
	\end{table}
	
	\subsection{Replicative lifespan determination}
	Cells are passaged at 80\% confluency and cumulative population doublings (CPD) are recorded. The primary endpoint is the CPD at growth arrest. BCLM‑optimised cells are predicted to reach \(1.5\)–\(2.0\times\) the CPD of WT‑OE controls.
	
	\section{Cross‑Species Controlled Coordination Regression (CCR)}
	\label{sec:CCR}
	The reversibility of the ageing phenotype is tested in each species by transiently suppressing one optimised gene per layer.
	
	\begin{table}[H]
		\centering
		\caption{CCR interventions for the five species.}
		\label{tab:CCR}
		\begin{tabular}{@{}llll@{}}
			\toprule
			\textbf{Layer} & \textbf{Target gene} & \textbf{Method} & \textbf{Expected effect} \\
			\midrule
			1 & BRCA1\_LL & Dox‑inducible shRNA & \(\REL \to 0.62\); mutation rate rises by factor \(1.86\) \\
			2 & NDUFS1\_LL & CRISPR base‑edit (I182F) & \(\REL \to 0.55\); ROS rises by factor \(1.80\) \\
			3 & HSPA1A\_LL & Tet‑CRISPRi & \(\REL \to 0.52\); aggregates rise by factor \(1.78\) \\
			4 & ITGB3\_LL & siRNA & \(\REL \to 0.57\); migration speed falls by factor \(1.20\) \\
			\bottomrule
		\end{tabular}
	\end{table}
	
	After 72 h of suppression, biomarkers must shift toward the wild‑type phenotype. Following washout or siRNA dilution, cells should recover the Long‑Life state within 48–72 h. The quantitative prediction for each biomarker is derived directly from the \(\REL\) values in Table~\ref{tab:targets} and the universal error law.
	
	\section{Falsifiability}
	\label{sec:falsifiability}
	All predictions in Table~\ref{tab:predictions} have been preregistered in the YPSDC repository. The YUCT coordination theory of ageing will be considered falsified if, after three independent repeats, the experimental outcomes fall outside the stated ranges for four out of five species.
	
	\begin{table}[H]
		\centering
		\caption{Pregistered predictions (identical for all five species).}
		\label{tab:predictions}
		\begin{tabular}{@{}llll@{}}
			\toprule
			\textbf{Prediction} & \textbf{Assay} & \textbf{Expected (LL/WT)} & \textbf{Falsification if} \\
			\midrule
			Mutation rate       & HPRT           & \(0.86 \pm 0.05\) & LL/WT \(> 0.95\) \\
			ROS                 & MitoSOX        & \(0.80 \pm 0.05\) & LL/WT \(> 0.90\) \\
			Aggregates          & HTT‑Q72 foci   & \(0.78 \pm 0.05\) & LL/WT \(> 0.90\) \\
			Migration speed     & Scratch assay  & \(1.20 \pm 0.05\) & LL/WT \(< 1.05\) \\
			Replicative lifespan & Cumulative PD  & \(1.5\)–\(2.0\)   & LL/WT \(< 1.3\) \\
			\bottomrule
		\end{tabular}
	\end{table}
	
	\section{Ethical Considerations}
	\label{sec:ethics}
	This protocol uses exclusively commercially available cell lines. No live animals are required. The genetic constructs are delivered \textit{in vitro} via standard transfection reagents; \textit{in vivo} delivery methods are not described. Researchers are urged to consult Appendix~\(\Sigma\) before extending this work beyond cultured cells.
	
	\section{Conclusion}
	\label{sec:conclusion}
	The BCLM‑LL‑MPS protocol provides a rigorous, internally consistent framework to test the YUCT coordination theory of ageing across five evolutionarily distinct mammalian species. All predictions are computed from a single set of universal constants and species‑specific codon tables; no free parameters are adjusted \textit{post hoc}. A successful outcome would establish \(K_{\mathrm{eff}}\) as the first cross‑species, quantitative metric of biological coordination and its restoration as a rational anti‑ageing strategy.
	
	\vspace{0.5cm}
	\noindent\textbf{Data and code:} All orthologous sequences, codon‑usage tables, computed \(\REL\) values, and the preregistered predictions are at \url{https://github.com/Alexey-Yakushev-YUCT/YPSDC}.
	
	\begin{thebibliography}{99}
		\bibitem{BCLM} A.\,V.\,Yakushev, ``Appendix BCLM: Biological Coordination Lagrangian,'' in \emph{YUCT}, Zenodo (2026).
		\bibitem{BCLM_LL} A.\,V.\,Yakushev, ``Appendix BCLM‑LL: Experimental Protocol … Human Cardiomyocytes,'' in \emph{YUCT}, Zenodo (2026).
		\bibitem{AppendixL} A.\,V.\,Yakushev, ``Appendix L: Fractal Coordination Error Scaling,'' in \emph{YUCT}, Zenodo (2026).
		\bibitem{AppendixP} A.\,V.\,Yakushev, ``Appendix P: Temperature Dependence of Gravity,'' in \emph{YUCT}, Zenodo (2026).
		\bibitem{AppendixSigma} A.\,V.\,Yakushev, ``Appendix \(\Sigma\): Ethical Imperatives and Governance,'' in \emph{YUCT}, Zenodo (2026).
		\bibitem{Bjedov2021} I.~Bjedov \emph{et al.}, ``A single hyper‑accurate mutation in the ribosome extends lifespan in yeast, worms, and flies,'' \emph{Cell Metabolism}, 33, 1054–1070 (2021).
		\bibitem{Kerekes2025} K.~Kerekes \emph{et al.}, ``Codon optimisation and evolutionary pressure in long‑lived mammals,'' \emph{bioRxiv} (2025).
		\bibitem{Morrow2004} G.~Morrow \emph{et al.}, ``Overexpression of … Hsp22 extends \textit{Drosophila} lifespan,'' \emph{FASEB J.}, 18, 598–599 (2004).
		\bibitem{Walker2003} G.\,A.~Walker, G.\,J.~Lithgow, ``Lifespan extension in \textit{C. elegans} by a molecular chaperone,'' \emph{Aging Cell}, 2, 131–139 (2003).
		\bibitem{Takeda2019} M.~Takeda \emph{et al.}, ``Downregulation of \(\alpha\)5 integrin induces cellular senescence,'' \emph{FEBS Lett.}, 593, 1284–1293 (2019).
		\bibitem{codon_usage_db} Y.\,Nakamura \emph{et al.}, Codon usage tabulated … \url{https://www.kazusa.or.jp/codon/}.
		\bibitem{YPSDC_repo} A.\,V.\,Yakushev, YPSDC Repository, \url{https://github.com/Alexey-Yakushev-YUCT/YPSDC}.
	\end{thebibliography}
	
\end{document}